The ten properties are:
\begin{itemize}[leftmargin=*,itemsep=1pt]
\item \textbf{P1:} The evaluated model acts by emitting tool calls. A model that only judges as a monitor does not meet this property.
\item \textbf{P2:} The unsafe outcome is scored from a change in the environment's state.
\item \textbf{P3:} Attempted actions are reported separately from executed actions.
\item \textbf{P4:} A signed or cryptographic artifact varies as an independent variable.
\item \textbf{P5:} The evaluation tests signed-but-invalid forms, such as replayed artifacts, wrong binding, wrong keys or artifacts outside their validity window.
\item \textbf{P6:} The delivery channel of the evidence or injection is analyzed as a factor.
\item \textbf{P7:} A pre-registration declares the prompt-wording factor.
\item \textbf{P8:} The evaluation includes a code-enforced check at the point of action.
\item \textbf{P9:} The primary outcome is scored deterministically, without a judge model.
\item \textbf{P10:} The acting agent itself receives the varied clearance or authorization evidence.
\end{itemize}

Three entries for this study are partial. For P5, forged signatures and unregistered signing keys appear in verifier unit tests but are not episode variants. For P7, the pre-registration is a dated document committed by hash and was not filed with an external registry. For P8, \arm{A2} serves as an enforcement baseline: its code excludes unsafe executions by construction. $^{\dagger}$FinVault was withdrawn by its authors in arXiv v2 on 30 July 2026. Its table entries describe v1.
